\section{Related work}
\label{sec:related-work}

The published work that bears on the proposed method falls into five groups: the estimators of VaR and ES that the wrapper surrounds, the back-tests and scores that judge it, the evidence on how bank models behave in stress, the conformal methods that supply the update, and the recalibrations of risk forecasts that already exist.

\subsection{Estimators of value at risk and expected shortfall}
\label{sec:rw-estimators}

Historical simulation takes the empirical quantile of a fixed window of past returns and needs no distributional assumption; it reacts slowly to a rise in volatility, because a crisis day enters a window of 250 days with the same weight as a day of the calm year, and it cannot resolve a tail level below the reciprocal of the window \citep{pritsker2006}. The age-weighted variant gives recent days more weight through a geometric decay \citep{boudoukh1998}, and volatility rescaling brings the history to the current volatility \citep{hullwhite1998}. Parametric models fit a conditional scale: the generalised autoregressive conditional heteroskedasticity model of \citet{bollerslev1986} and its asymmetric forms, the model of \citet{gjr1993} and the exponential model of \citet{nelson1991}, whose innovations are Student-$t$ in the ladder used here. Filtered historical simulation draws the standardised residuals of such a filter from their own empirical distribution \citep{barone1999}, and the conditional extreme-value model of \citet{mcneilfrey2000} fits a generalised Pareto law to the tail of the residuals. The conditional autoregressive quantile models of \citet{engle2004} regress the quantile on its own past and on the absolute return and give no ES; \citet{taylor2019} adds a second step that estimates the ES from the gap below the quantile, and \citet{pattonziegelchen2019} fit the VaR and ES pair with dynamic recursions driven by the score of the joint loss. \citet{kuester2006} compare alternative strategies for VaR prediction; all of these estimators assume that the parameters estimated on the past hold on the day to come. The wrapper does not compete with any of them: it asks each one for its forecast at another level.

\subsection{Back-tests and scores for the pair}
\label{sec:rw-backtests}

The proportion-of-failures test of \citet{kupiec1995} compares the frequency of exceedances with the tail level and has low power at the sample sizes of a one-year window. \citet{christoffersen1998} adds a test of the independence of the exceedances against a first-order Markov alternative and combines the two into a conditional coverage test; \citet{christoffersenpelletier2004} replace the Markov alternative by the durations between exceedances, and \citet{berkowitz2011} evaluate VaR models with the profit and loss of the individual trading desks, which gives the tests more observations than the aggregate bank does. \citet{escanciano2012} show that the estimation error of a historical-simulation VaR distorts the size of these tests. Tests for the ES use the size of the losses beyond the VaR: the cumulative-violation process of \citet{duescanciano2017}, the multinomial tests of \citet{kratz2018}, the regression tests of \citet{bayer2022}, and the e-process backtests of \citet{wang2022}. The scoring side starts from the fact that a quantile is elicitable and the ES is not, and that the pair is jointly elicitable \citep{gneiting2011,fisslerziegelgneiting2016,fissler2016,ziegel2016}; the loss $\mathrm{FZ0}$ of \citet{pattonziegelchen2019} is a member of the family of consistent scores for the pair and the one used here. Comparisons of models under such a score use the test of \citet{diebold1995} and the model confidence set of \citet{hansen2011}; \citet{noldeziegel2017} study how the scoring view and the supervisory back-test differ. The back-tests used by the supervisor are specified in the Basel framework \citep{bcbsmar32,bcbsmar33}.

\subsection{Bank practice in stress}
\label{sec:rw-practice}

\citet{berkowitz2002} assess the accuracy of the VaR models of commercial banks against their realised trading results, \citet{perignon2008} ask whether banks overstate their VaR, and \citet{perignon2010} study the level and the quality of the VaR that banks disclose. These studies of reported VaR concern its average level in calm periods and its disclosure; the behaviour at the onset of a crisis, where the exceedances cluster, is what the stress windows of this manuscript examine. The supervisory record of that behaviour is the count of exceptions of the 99\% VaR over 250 trading days. The framework of \citet{bcbs1996} sets the green zone where the cumulative binomial probability of the count, at the reporting frequency of 1\%, stays below 95\%, which is four exceptions or fewer (the exact probability of at most four is 89.2\% and of at most five 95.9\%), and the red zone where it reaches 99.99\%, which is ten exceptions or more. The multiplier on the VaR-based capital requirement follows the zone; in the United States implementation it is 3.00 for four exceptions or fewer, 3.40 at five, rising to 3.85 at nine, and 4.00 from ten \citep{cfr217204}. The remedy therefore acts on capital after the count has been reached, and it is blind to the timing of the exceptions: ten exceptions spread over the year and ten in one month receive the same multiplier. The coherence of the ES as a risk measure motivates the supervisor's move from the VaR to the ES \citep{artzner1999,acerbitasche2002}.

\subsection{Conformal methods for online data}
\label{sec:rw-conformal}

Conformal prediction builds a prediction set with a coverage guarantee that holds without a distributional assumption beyond exchangeability \citep{vovk2005}; the split form \citep{lei2018} calibrates the width on a held-out sample, and the quantile form \citep{romano2019} lets the width depend on the covariates. A tutorial treatment is in \citet{angelopoulosbates2023}. Exchangeability fails for time series, and three lines of work address it. The first reweights the calibration scores to account for a shift \citep{tibshirani2019,barber2023}, which needs knowledge of the shift. The second builds the interval from the residuals of an ensemble of models trained on blocks of the series \citep{xu2021}. The third adapts the miscoverage level online: the adaptive conformal inference of \citet{gibbs2021} moves the level by a constant multiple of the miscoverage indicator minus the target and guarantees the long-run coverage for any data sequence; \citet{zaffran2022} aggregate the levels obtained from a grid of steps, \citet{gibbs2024} adapt the step itself, \citet{bhatnagar2023} obtain coverage on every interval of the sequence with a strongly adaptive learner, \citet{angelopoulos2023} add an integrator and a forecaster of the scores to the update, forming a control loop of three terms, and \citet{angelopoulos2024} give decaying steps with a convergence rate. These results concern intervals; the guarantee is on the frequency of the miscoverage and not on its independence.

\subsection{Recalibration of risk forecasts and the position of the proposed method}
\label{sec:rw-position}

The closest published works apply conformal updates to the VaR. \citet{fantazzini2024} run four adaptive conformal algorithms on the log-returns of four thousand crypto-assets and compare them with GARCH models and models of the daily range; they report that the conformal VaR holds its level where the parametric models suffer structural breaks and numerical failures, and that the aggregated variants are conservative at higher levels. \citet{aich2025} study two rolling conformal procedures for intervals of returns on thirteen series. Neither wraps a model that is already in a bank's production, and neither reports the Kupiec, independence and conditional coverage tests together with the joint score of the VaR and ES for each model of a ladder in the stress windows of the supervisor. The proposed method differs from these in the object it wraps (a full VaR and ES model of the ladder, asked at another level, with its own ES kept), in the judgement it is subjected to (the supervisory tests and the pair score, against the unwrapped model and against two recalibrations that act on the finished VaR, run on the same forecasts) and in the account of its limits (the identity that shows the guarantee is vacuous on a window of a year, the band-limited repair, and the effect of the band on the hit rate). The step update itself is the one of \citet{gibbs2021}; the manuscript claims no change of it. The connection to quantile regression \citep{koenker1978} is that the update is a stochastic-approximation step \citep{robbins1951} on the expected indicator, which the theory section uses.
